%! Inverse Trigonometric Functions — Unit 3, Lesson 3.9 — Homework
\documentclass[10pt]{article}
\usepackage{precalculus-article}
\usepackage{precalculus-boxes}
\newcommand{\TallMath}[1]{$\displaystyle #1\rule[-1.4em]{0pt}{3.2em}$}

\begin{document}

\pageheader{Unit 3, Lesson 3.9}{Homework}
\namedateperiod

\vspace{0.10in}

\begin{notesbox}{Part 1 --- Evaluate}

Evaluate each expression.  Give the exact answer in radians.

\renewcommand{\arraystretch}{2.0}
\begin{tabularx}{\linewidth}{@{} X X X @{}}
  \textbf{1.} $\arcsin\!\left(\dfrac{\sqrt{2}}{2}\right) = $ \blank{2cm}
  &
  \textbf{2.} $\arccos\!\left(-\dfrac{1}{2}\right) = $ \blank{2cm}
  &
  \textbf{3.} $\arctan\!\left(\dfrac{\sqrt{3}}{3}\right) = $ \blank{2cm} \\[4pt]
  \textbf{4.} $\sin^{-1}(0) = $ \blank{2cm}
  &
  \textbf{5.} $\cos^{-1}(1) = $ \blank{2cm}
  &
  \textbf{6.} $\tan^{-1}\!\left(-\dfrac{\sqrt{3}}{3}\right) = $ \blank{2cm}
\end{tabularx}

\end{notesbox}

\vspace{0.08in}

\begin{notesbox}{Part 2 --- Domain and Range}

\textbf{7.} State the domain and range of each function.

\renewcommand{\arraystretch}{1.9}
\begin{tabular}{|l|c|c|}
\hline
\textbf{Function} & \textbf{Domain} & \textbf{Range} \\
\hline
$y = \arcsin x$    & \blank{4cm} & \blank{4cm} \\
\hline
$y = \arccos x$    & \blank{4cm} & \blank{4cm} \\
\hline
$y = \arctan x$    & \blank{4cm} & \blank{4cm} \\
\hline
\end{tabular}

\end{notesbox}

\vspace{0.08in}

\begin{notesbox}{Part 3 --- AP-Style Multiple Choice}

\textbf{8.} Which of the following is the value of $\arccos\!\left(-\dfrac{\sqrt{2}}{2}\right)$?

\begin{tabularx}{\linewidth}{@{} X X X X X @{}}
(A) $-\dfrac{3\pi}{4}$ &
(B) $-\dfrac{\pi}{4}$ &
(C) $\dfrac{\pi}{4}$ &
(D) $\dfrac{3\pi}{4}$ &
(E) $\dfrac{5\pi}{4}$
\end{tabularx}

\vspace{0.10in}

\textbf{9.} Let $f(x) = \arctan(x)$.  Which statement about $f$ is true?

\begin{enumerate}[leftmargin=2em, label=(\Alph*), itemsep=3pt]
    \item The range of $f$ is all real numbers.
    \item $f(x)$ is undefined for $x < -1$.
    \item The range of $f$ is $\left(-\dfrac{\pi}{2},\,\dfrac{\pi}{2}\right)$.
    \item $f(-x) = -f(x)$ is false for all $x \ne 0$.
\end{enumerate}

\end{notesbox}

\vspace{0.08in}

\begin{extensionbox}{Extension --- Composition of Trig and Inverse Trig}

\textbf{10.} Simplify $\sin(\arctan x)$ for $x > 0$.  Express your answer
without any trig or inverse trig functions.

\writelines{4}

\vspace{0.06in}

\textbf{Preview Lesson 3.10:} Consider the equation $\cos\theta = \tfrac{1}{2}$.
How many solutions exist on $[0, 2\pi)$?  List them.  How would you use
$\arccos$ to find one solution?  We will explore this in the next lesson.

\writelines{2}

\end{extensionbox}

\end{document}
